\chapter{1998年全国卷（理）}
\section{单选题}

\begin{problem}
\(\sin 600^{\circ}\) 的值是（\quad）

\choices
    {\(\frac{1}{2}\)}
    {\(-\frac{1}{2}\)}
    {\(\frac{\sqrt{3}}{2}\)}
    {\(-\frac{\sqrt{3}}{2}\)}
\end{problem}

\begin{answer}
D
\end{answer}

\begin{solution}
因为 \(600^{\circ}=360^{\circ}+240^{\circ}\)，所以
\(\sin 600^{\circ}=\sin 240^{\circ}=-\sin 60^{\circ}=-\frac{\sqrt{3}}{2}\)，故选 D.
\end{solution}

\begin{problem}
函数 \(y=a^{|x|}\)（\(a>1\)）的图象是（\quad）

\choices
    {\choicebitmap[width=0.15\paperwidth]{img/1998/national_science/q2_optA_nolabel.png}}
    {\choicebitmap[width=0.15\paperwidth]{img/1998/national_science/q2_optB_nolabel.png}}
    {\choicebitmap[width=0.15\paperwidth]{img/1998/national_science/q2_optC_nolabel.png}}
    {\choicebitmap[width=0.15\paperwidth]{img/1998/national_science/q2_optD_nolabel.png}}
\end{problem}

\begin{answer}
B
\end{answer}

\begin{solution}
由于 \(a>1\) 且 \(|x|\geq0\)，有 \(a^{|x|}\geq1\)，等号仅在 \(x=0\) 时成立；又 \(y(-x)=y(x)\)，所以图象关于 \(y\) 轴对称，并在点 \((0,1)\) 处取得最小值. 随着 \(|x|\) 增大，\(a^{|x|}\) 增大，因此图象为选项 B.
\end{solution}

\begin{problem}
曲线的极坐标方程 \(\rho = 4\sin \theta \) 化成直角坐标方程为（\quad）

\choices
    {\(x^{2} + (y + 2)^{2} = 4\)}
    {\(x^{2} + (y - 2)^{2} = 4\)}
    {\((x - 2)^{2} + y^{2} = 4\)}
    {\((x + 2)^{2} + y^{2} = 4\)}
\end{problem}

\begin{answer}
B
\end{answer}

\begin{solution}
由极坐标与直角坐标的关系 \(\rho^{2}=x^{2}+y^{2}\)，\(\rho\sin\theta=y\)，原方程化为
\(x^{2}+y^{2}=4y\)，即 \(x^{2}+(y-2)^{2}=4\)，故选 B.
\end{solution}

\begin{problem}
两条直线 \(A_{1}x + B_{1}y + C_{1} = 0\)，\(A_{2}x + B_{2}y + C_{2} = 0\) 垂直的充要条件是（\quad）

\choices
    {\(A_{1}A_{2} + B_{1}B_{2} = 0\)}
    {\(A_{1}A_{2} - B_{1}B_{2} = 0\)}
    {\(\frac{A_1A_2}{B_1B_2} = -1\)}
    {\(\frac{B_1B_2}{A_1A_2} = 1\)}
\end{problem}

\begin{answer}
A
\end{answer}

\begin{solution}
两条直线的方向向量可分别取 \((B_{1},-A_{1})\) 与 \((B_{2},-A_{2})\). 两直线垂直当且仅当这两个方向向量的内积为零，即
\(B_{1}B_{2}+A_{1}A_{2}=0\)，所以充要条件为 \(A_{1}A_{2}+B_{1}B_{2}=0\)，故选 A.
\end{solution}

\begin{problem}
函数 \(f(x)=\frac{1}{x}\)（\(x\neq 0\)）的反函数 \(f^{-1}(x)=\)（\quad）

\choices
    {\(x\)（\(x\neq 0\)）}
    {\(\frac{1}{x}\)（\(x\neq 0\)）}
    {\(-x\)（\(x\neq 0\)）}
    {\(-\frac{1}{x}\)（\(x\neq 0\)）}
\end{problem}

\begin{answer}
B
\end{answer}

\begin{solution}
设 \(y=f(x)=\frac{1}{x}\)，则 \(x=\frac{1}{y}\)，且 \(x\neq0\)、\(y\neq0\). 交换 \(x\)、\(y\) 得反函数 \(y=\frac{1}{x}\)，所以 \(f^{-1}(x)=\frac{1}{x}\)（\(x\neq0\)），故选 B.
\end{solution}

\begin{problem}
已知点 \(P(\sin \alpha - \cos \alpha, \tan \alpha)\) 在第一象限，则在 \([0, 2\pi]\) 内 \(\alpha\) 的取值范围是（\quad）

\choices
    {\(\left(\frac{\pi}{2}, \frac{3\pi}{4}\right) \cup \left(\pi, \frac{5\pi}{4}\right)\)}
    {\(\left(\frac{\pi}{4}, \frac{\pi}{2}\right) \cup \left(\pi, \frac{5\pi}{4}\right)\)}
    {\(\left(\frac{\pi}{2}, \frac{3\pi}{4}\right) \cup \left(\frac{5\pi}{4}, \frac{3\pi}{2}\right)\)}
    {\(\left(\frac{\pi}{4}, \frac{\pi}{2}\right) \cup \left(\frac{3\pi}{4}, \pi\right)\)}
\end{problem}

\begin{answer}
B
\end{answer}

\begin{solution}
点 \(P\) 在第一象限等价于
\(\tan\alpha>0\)，且 \(\sin\alpha-\cos\alpha>0\). 在 \([0,2\pi]\) 内，\(\tan\alpha>0\) 给出
\(\alpha\in(0,\frac{\pi}{2})\cup(\pi,\frac{3\pi}{2})\).

在第一象限内，\(\sin\alpha>\cos\alpha\) 等价于 \(\alpha>\frac{\pi}{4}\)，得到 \(\alpha\in(\frac{\pi}{4},\frac{\pi}{2})\). 在第三象限内，利用
\(\sin\alpha-\cos\alpha=\sqrt{2}\sin(\alpha-\frac{\pi}{4})\)，得到 \(\alpha\in(\pi,\frac{5\pi}{4})\). 因此
\(\alpha\in(\frac{\pi}{4},\frac{\pi}{2})\cup(\pi,\frac{5\pi}{4})\)，故选 B.
\end{solution}

\begin{problem}
已知圆锥的全面积是底面积的 \(3\) 倍，那么该圆锥的侧面展开图扇形的圆心角为（\quad）

\choices
    {\(120^{\circ}\)}
    {\(150^{\circ}\)}
    {\(180^{\circ}\)}
    {\(240^{\circ}\)}
\end{problem}

\begin{answer}
C
\end{answer}

\begin{solution}
设圆锥底面半径为 \(r\)，母线长为 \(l\). 由全面积等于底面积的 \(3\) 倍，得侧面积等于底面积的 \(2\) 倍，即
\(\pi rl=2\pi r^{2}\)，所以 \(l=2r\).

设侧面展开图扇形的圆心角为 \(\theta\)（弧度），扇形弧长等于底面周长，故 \(\theta l=2\pi r\). 代入 \(l=2r\)，得 \(\theta=\pi=180^{\circ}\)，故选 C.
\end{solution}

\begin{problem}
复数 \(-\i\) 的一个立方根是 \(\i\)，它的另外两个立方根是（\quad）

\choices
    {\(\frac{\sqrt{3}}{2} \pm \frac{1}{2}\i\)}
    {\(-\frac{\sqrt{3}}{2} \pm \frac{1}{2}\i\)}
    {\(\pm \frac{\sqrt{3}}{2} + \frac{1}{2}\i\)}
    {\(\pm \frac{\sqrt{3}}{2} - \frac{1}{2}\i\)}
\end{problem}

\begin{answer}
D
\end{answer}

\begin{solution}
设 \(\omega=-\frac{1}{2}+\frac{\sqrt{3}}{2}\i\) 为三次单位根，则已知立方根 \(\i\) 的另外两个立方根为 \(\i\omega\) 和 \(\i\omega^{2}\). 分别计算得
\(\i\omega=-\frac{\sqrt{3}}{2}-\frac{1}{2}\i\)，\(\i\omega^{2}=\frac{\sqrt{3}}{2}-\frac{1}{2}\i\). 因此另外两个根为 \(\pm\frac{\sqrt{3}}{2}-\frac{1}{2}\i\)，故选 D.
\end{solution}

\begin{problem}
如果棱台的两底面积分别是 \(S\)，\(S'\)，中截面的面积是 \(S_0\). 那么（\quad）

\choices
    {\(2\sqrt{S_0} = \sqrt{S} +\sqrt{S'}\)}
    {\(S_0 = \sqrt{SS'}\)}
    {\(2S_{0} = S + S^{\prime}\)}
    {\(S_0^2 = 2SS'\)}
\end{problem}

\begin{answer}
A
\end{answer}

\begin{solution}
棱台中截面各对应线段的长度等于两底面对应线段长度的平均值. 因而中截面的面积平方根等于两底面积平方根的平均值，即
\(\sqrt{S_{0}}=\frac{\sqrt{S}+\sqrt{S'}}{2}\). 所以
\(2\sqrt{S_{0}}=\sqrt{S}+\sqrt{S'}\)，故选 A.
\end{solution}

\begin{problem}
向高为 \(H\) 的水瓶中注水，注满为止，如果注水量 \(V\) 与水深 \(h\) 的函数关系的图象如下图所示，那么水瓶的形状是（\quad）
\bitmapfigure[max width=0.30\linewidth]{img/1998/national_science/q10_fig1.png}

\choices
    {\choicebitmap[width=0.15\paperwidth]{img/1998/national_science/q10_optA_nolabel.png}}
    {\choicebitmap[width=0.15\paperwidth]{img/1998/national_science/q10_optB_nolabel.png}}
    {\choicebitmap[width=0.15\paperwidth]{img/1998/national_science/q10_optC_nolabel.png}}
    {\choicebitmap[width=0.15\paperwidth]{img/1998/national_science/q10_optD_nolabel.png}}
\end{problem}

\begin{answer}
B
\end{answer}

\begin{solution}
从图象可知，随着水深 \(h\) 增加，相同水深增量所对应的注水量增量逐渐减小. 这说明水瓶在各高度处的水平截面积自下而上逐渐减小，即瓶底较宽、瓶口较窄. 四个图形中符合这一特征的是选项 B.
\end{solution}

\begin{problem}
\(3\) 名医生和 \(6\) 名护士被分配到 \(3\) 所学校为学生体检，每校分配 \(1\) 名医生和 \(2\) 名护士，不同的分配方法共有（\quad）

\choices
    {\(90\) 种}
    {\(180\) 种}
    {\(270\) 种}
    {\(540\) 种}
\end{problem}

\begin{answer}
D
\end{answer}

\begin{solution}
先把 \(3\) 名不同的医生分配到 \(3\) 所不同的学校，有 \(3!\) 种方法. 再把 \(6\) 名不同的护士按学校依次分成三组，每组 \(2\) 人，有
\(\mathrm C_{6}^{2}\mathrm C_{4}^{2}\mathrm C_{2}^{2}=90\) 种方法. 因此总方法数为
\(3!\mathrm C_{6}^{2}\mathrm C_{4}^{2}\mathrm C_{2}^{2}=6\times90=540\)，故选 D.
\end{solution}

\begin{problem}
椭圆 \(\frac{x^2}{12} + \frac{y^2}{3} = 1\) 的焦点为 \(F_{1}\) 和 \(F_{2}\). 点 \(P\) 在椭圆上，如果线段 \(PF_{1}\) 的中点在 \(y\) 轴上，那么 \(|PF_{1}|\) 是 \(|PF_{2}|\) 的（\quad）

\choices
    {\(7\) 倍}
    {\(5\) 倍}
    {\(4\) 倍}
    {\(3\) 倍}
\end{problem}

\begin{answer}
A
\end{answer}

\begin{solution}
椭圆的半长轴平方为 \(12\)，半短轴平方为 \(3\)，所以焦半距满足 \(c^{2}=12-3=9\)，可取 \(F_{1}=(-3,0)\)、\(F_{2}=(3,0)\). 设 \(P=(x_{P},y_{P})\)，由 \(PF_{1}\) 的中点在 \(y\) 轴上，得
\(\frac{x_{P}-3}{2}=0\)，所以 \(x_{P}=3\). 代入椭圆方程得
\(\frac{9}{12}+\frac{y_{P}^{2}}{3}=1\)，从而 \(y_{P}^{2}=\frac{3}{4}\). 于是
\(|PF_{1}|^{2}=(3+3)^{2}+\frac{3}{4}=\frac{147}{4}\)，而 \(|PF_{2}|^{2}=\frac{3}{4}\)，故
\(\frac{|PF_{1}|}{|PF_{2}|}=\sqrt{\frac{147}{3}}=7\)，所以选 A.
\end{solution}

\begin{problem}
球面上有 \(3\) 个点，其中任意两点的球面距离都等于大圆周长的 \(\frac{1}{6}\)，经过这 \(3\) 个点的小圆的周长为 \(4\pi\)，那么这个球的半径为（\quad）

\choices
    {\(4\sqrt{3}\)}
    {\(2\sqrt{3}\)}
    {\(2\)}
    {\(\sqrt{3}\)}
\end{problem}

\begin{answer}
B
\end{answer}

\begin{solution}
设球的半径为 \(R\)，经过三个点的小圆半径为 \(r\). 由小圆周长为 \(4\pi\)，得 \(2\pi r=4\pi\)，所以 \(r=2\).

任意两点的球面距离为大圆周长的 \(\frac{1}{6}\)，故它们在球心所对的圆心角为 \(60^{\circ}\)，两点间的弦长为
\(2R\sin30^{\circ}=R\). 三个点在同一小圆上且两两弦长相等，因此构成内接于该小圆的正三角形，其边长为 \(\sqrt{3}r=2\sqrt{3}\). 于是 \(R=2\sqrt{3}\)，故选 B.
\end{solution}

\begin{problem}
一个直角三角形三内角的正弦值成等比数列，其最小内角为（\quad）

\choices
    {\(\arccos \frac{\sqrt{5} - 1}{2}\)}
    {\(\arcsin \frac{\sqrt{5} - 1}{2}\)}
    {\(\arccos \frac{1 - \sqrt{5}}{2}\)}
    {\(\arcsin \frac{1 - \sqrt{5}}{2}\)}
\end{problem}

\begin{answer}
B
\end{answer}

\begin{solution}
设直角三角形的两个锐角中较小者为 \(\alpha\)，另一个为 \(\beta\)，则 \(\alpha+\beta=\frac{\pi}{2}\)，且
\(0<\alpha\leq\beta<\frac{\pi}{2}\). 三个内角的正弦值依次为 \(\sin\alpha\)、\(\sin\beta\)、\(1\)，它们成等比数列，所以
\(\sin^{2}\beta=\sin\alpha\). 又 \(\sin\beta=\cos\alpha\)，故 \(\cos^{2}\alpha=\sin\alpha\). 令 \(u=\sin\alpha\)，则
\(1-u^{2}=u\)，即 \(u^{2}+u-1=0\). 由于 \(0<u<1\)，得 \(u=\frac{\sqrt{5}-1}{2}\). 因此
\(\alpha=\arcsin\frac{\sqrt{5}-1}{2}\)，故选 B.
\end{solution}

\begin{problem}
在等比数列 \(\{a_{n}\}\) 中，\(a_{1}>1\)，且前 \(n\) 项和 \(S_{n}\) 满足 \(\displaystyle\lim_{n\to\infty} S_{n}=\frac{1}{a_{1}}\)，那么 \(a_{1}\) 的取值范围是（\quad）

\choices
    {\((1, +\infty)\)}
    {\((1,4)\)}
    {\((1,2)\)}
    {\((1,\sqrt{2})\)}
\end{problem}

\begin{answer}
D
\end{answer}

\begin{solution}
设等比数列的公比为 \(q\). 前 \(n\) 项和有极限，故 \(|q|<1\)，并且
\(\lim_{n\to\infty}S_{n}=\frac{a_{1}}{1-q}=\frac{1}{a_{1}}\). 由此得
\(q=1-a_{1}^{2}\). 结合 \(-1<q<1\)，可得 \(0<a_{1}^{2}<2\). 再由 \(a_{1}>1\)，得到
\(1<a_{1}<\sqrt{2}\)，故选 D.
\end{solution}

\section{填空题}

\begin{problem}
设圆过双曲线 \(\frac{x^{2}}{9} - \frac{y^{2}}{16} = 1\) 的一个顶点和一个焦点，圆心在此双曲线上，则圆心到双曲线中心的距离是\fillinblank{}.
\end{problem}

\begin{answer}
\(\frac{16}{3}\)
\end{answer}

\begin{solution}
双曲线的半实轴长为 \(3\)，半虚轴长为 \(4\)，所以焦半距满足
\(c=\sqrt{3^{2}+4^{2}}=5\). 设所取顶点为 \(V=(3\varepsilon,0)\)，所取焦点为 \(F=(5\delta,0)\)，其中 \(\varepsilon\)，\(\delta\in\{1,-1\}\)，圆心为 \(O=(u,v)\).

由于 \(O\) 到 \(V\)、\(F\) 的距离相等，\(O\) 在 \(VF\) 的垂直平分线上，从而
\(u=\frac{3\varepsilon+5\delta}{2}\). 若 \(\varepsilon\neq\delta\)，则 \(u=\pm1\)，但双曲线上的点满足 \(|u|\geq3\)，不可能. 因此 \(\varepsilon=\delta\)，从而 \(u=\pm4\).

将 \(u=\pm4\) 代入双曲线方程，得
\(\frac{16}{9}-\frac{v^{2}}{16}=1\)，所以 \(v^{2}=\frac{112}{9}\). 于是
\(O\) 到双曲线中心的距离满足
\(\lvert O\rvert^{2}=u^{2}+v^{2}=16+\frac{112}{9}=\frac{256}{9}\)，故 \(\lvert O\rvert=\frac{16}{3}\).
\end{solution}

\begin{problem}
\((x + 2)^{10}(x^{2} - 1)\) 的展开式中 \(x^{10}\) 的系数为\fillinblank{}. （用数字作答）
\end{problem}

\begin{answer}
\(179\)
\end{answer}

\begin{solution}
要得到 \(x^{10}\) 项，第一部分取 \(x^{8}\) 项并与第二部分的 \(x^{2}\) 相乘，或第一部分取 \(x^{10}\) 项并与第二部分的 \(-1\) 相乘. 因此所求系数为
\(\mathrm C_{10}^{8}2^{2}-1=45\times4-1=179\).
\end{solution}

\begin{problem}
在直四棱柱 \(A_{1}B_{1}C_{1}D_{1}-ABCD\) 中，当底面四边形 \(ABCD\) 满足条件\fillinblank{} 时，有 \(A_{1}C \perp B_{1}D_{1}\). （注：填上你认为正确的一种条件即可，不必考虑所有可能的情形）
\end{problem}

\begin{answer}
\(AC\perp BD\)
\end{answer}

\begin{solution}
设 \(\bs{w}=\overrightarrow{AA_{1}}\). 因为是直四棱柱，\(\bs{w}\) 垂直于底面 \(ABCD\)，所以 \(\bs{w}\cdot\overrightarrow{BD}=0\). 又有
\(\overrightarrow{A_{1}C}=\overrightarrow{AC}-\bs{w}\)，\(\overrightarrow{B_{1}D_{1}}=\overrightarrow{BD}\)，从而
\(\overrightarrow{A_{1}C}\cdot\overrightarrow{B_{1}D_{1}}=(\overrightarrow{AC}-\bs{w})\cdot\overrightarrow{BD}=\overrightarrow{AC}\cdot\overrightarrow{BD}\).

若底面两条对角线满足 \(AC\perp BD\)，则 \(\overrightarrow{AC}\cdot\overrightarrow{BD}=0\)，于是
\(\overrightarrow{A_{1}C}\cdot\overrightarrow{B_{1}D_{1}}=0\)，即 \(A_{1}C\perp B_{1}D_{1}\). 因此填 \(AC\perp BD\) 即可.
\end{solution}

\begin{problem}
关于函数 \(f(x)=4\sin \left(2x+\frac{\pi}{3}\right)\)（\(x\in\R\)），有下列命题：
\begin{circlelist}
\item 由 \(f(x_{1})=f(x_{2})=0\) 可得 \(x_{1}-x_{2}\) 必是 \(\pi\) 的整数倍；
\item 若 \(y=f(x)\) 的表达式可改写为 \(y=4\cos \left(2x-\frac{\pi}{6}\right)\)；
\item \(y=f(x)\) 的图象关于点 \(\left(-\frac{\pi}{6}, 0\right)\) 对称；
\item \(y=f(x)\) 的图象关于直线 \(x=-\frac{\pi}{6}\) 对称.
\end{circlelist}
其中正确的命题的序号是\fillinblank{}. （注：把你认为正确的命题的序号都填上）
\end{problem}

\begin{answer}
\(2\)，\(3\)
\end{answer}

\begin{solution}
命题 1 中，\(f(x)=0\) 等价于 \(2x+\frac{\pi}{3}=k\pi\)，其中 \(k\in\Z\)，所以零点为 \(x=\frac{k\pi}{2}-\frac{\pi}{6}\). 两个零点之差为 \(\frac{(k_{1}-k_{2})\pi}{2}\)，例如取相邻的两个零点时差为 \(\frac{\pi}{2}\)，不一定是 \(\pi\) 的整数倍，故命题 1 错误.

由 \(\sin u=\cos\left(\frac{\pi}{2}-u\right)\) 以及余弦函数的偶性，
\(4\sin\left(2x+\frac{\pi}{3}\right)=4\cos\left(\frac{\pi}{6}-2x\right)=4\cos\left(2x-\frac{\pi}{6}\right)\)，故命题 2 正确.

令 \(x_{0}=-\frac{\pi}{6}\). 对任意实数 \(t\)，有 \(f(x_{0}+t)=4\sin 2t\)，\(f(x_{0}-t)=-4\sin 2t\). 因此，曲线上点 \(\left(x_{0}+t,f(x_{0}+t)\right)\) 关于点 \(\left(x_{0},0\right)\) 的对称点为 \(\left(x_{0}-t,-f(x_{0}+t)\right)=\left(x_{0}-t,f(x_{0}-t)\right)\)，仍在曲线上，故命题 3 正确.

关于直线 \(x=x_{0}\) 对称要求 \(f(x_{0}+t)=f(x_{0}-t)\) 对任意 \(t\) 成立，但上面两式一般互为相反数. 例如取 \(t=\frac{\pi}{4}\)，两者分别为 \(4\) 和 \(-4\)，故命题 4 错误. 所以正确命题的序号为 \(2\)，\(3\).
\end{solution}

\section{解答题}

\begin{problem}
在 \(\triangle ABC\) 中，\(a\)，\(b\)，\(c\) 分别是角 \(A\)，\(B\)，\(C\) 的对边，设 \(a + c = 2b\)，\(A - C = \frac{\pi}{3}\)，求 \(\sin B\) 的值.
\end{problem}

\begin{solution}
由正弦定理以及 \(a+c=2b\)，得
\[
\sin A+\sin C=2\sin B
\]
由和差化积公式、\(A+B+C=\pi\) 及 \(A-C=\frac{\pi}{3}\)，得
\[
2\sin\frac{A+C}{2}\cos\frac{A-C}{2}=2\sin B
\]
\[
\sqrt{3}\cos\frac{B}{2}=2\sin B=4\sin\frac{B}{2}\cos\frac{B}{2}
\]
因为 \(0<B<\pi\)，所以 \(\cos\frac{B}{2}\ne0\)，从而
\[
\sin\frac{B}{2}=\frac{\sqrt{3}}{4}
\]
于是
\[
\cos\frac{B}{2}=\sqrt{1-\sin^{2}\frac{B}{2}}=\frac{\sqrt{13}}{4}
\]
因此
\[
\sin B=2\sin\frac{B}{2}\cos\frac{B}{2}=\frac{\sqrt{39}}{8}
\]
\end{solution}

\begin{problem}
如图，直线 \(l_{1}\) 和 \(l_{2}\) 相交于点 \(M\)，\(l_{1} \perp l_{2}\)，点 \(N \in l_{1}\)，以 \(A\)，\(B\) 为端点的曲线段 \(C\) 上的任一点到 \(l_{2}\) 的距离与到点 \(N\) 的距离相等. 若 \(\triangle AMN\) 为锐角三角形，\(|AM| = \sqrt{17}\)，\(|AN| = 3\)，且 \(|BN| = 6\). 建立适当的坐标系，求曲线段 \(C\) 的方程.

\bitmapfigure[max width=5.6cm]{img/1998/national_science/q21_fig1.png}
\end{problem}

\begin{solution}
取 \(l_{1}\) 为 \(x\) 轴，\(l_{2}\) 为 \(y\) 轴，交点 \(M\) 为原点，并令 \(N=(n,0)\)，其中 \(n>0\). 设曲线段 \(C\) 上任一点为 \(P=(x,y)\). 由题意，点 \(P\) 到 \(l_{2}\) 的距离为 \(\lvert x\rvert\)，到 \(N\) 的距离为 \(\sqrt{(x-n)^{2}+y^{2}}\)，故
\[
x^{2}=(x-n)^{2}+y^{2}
\]
即 \(y^{2}=2nx-n^{2}\).
设 \(A=(x_{A},y_{A})\). 因为 \(A\) 在曲线上，且 \(A\) 到 \(l_{2}\) 的距离等于到 \(N\) 的距离，所以 \(\lvert x_{A}\rvert=\lvert AN\rvert=3\)，即 \(x_{A}^{2}=9\). 若 \(x_{A}=-3\)，则 \(y_{A}^{2}=2n(-3)-n^{2}<0\)，不可能，故 \(x_{A}=3\). 再由 \(\lvert AM\rvert=\sqrt{17}\)，得 \(y_{A}^{2}=17-3^{2}=8\). 代入抛物线方程，得
\[
6n-n^{2}=8
\]
所以 \(n=2\) 或 \(n=4\). 当 \(n=2\) 时，\(\overrightarrow{NM}=(-2,0)\)，\(\overrightarrow{NA}=(1,y_{A})\)，两向量内积为 \(-2<0\)，故三角形 \(AMN\) 在 \(N\) 处为钝角，不合题意. 当 \(n=4\) 时，三边平方分别为
\(|AM|^{2}=17\)、\(|AN|^{2}=9\)、\(|MN|^{2}=16\)，且 \(17<9+16\)、\(16<17+9\)、\(9<17+16\)，所以三角形 \(AMN\) 为锐角三角形，故 \(n=4\).

对曲线段端点 \(B=(x_{B},y_{B})\)，同样有
\[
BN^{2}=(x_{B}-n)^{2}+y_{B}^{2}=x_{B}^{2}
\]
所以 \(x_{B}^{2}=BN^{2}=36\). 当 \(x_{B}=-6\) 时，\(y_{B}^{2}=8x_{B}-16<0\)，不可能，故 \(x_{B}=6\). 按图取上方分支，曲线段 \(C\) 的方程为
\(y^{2}=8(x-2)\)，\(3\leq x\leq6\)，\(y>0\)
\end{solution}

\begin{problem}
如图，为处理含有某种杂质的水污染，要制造一底宽为 \(2\text{ 米}\) 的无盖长方体沉淀箱，污水从 \(A\) 汇流入，经沉淀后从 \(B\) 汇流出. 设箱体的长度为 \(a\text{ 米}\)，高度为 \(b\text{ 米}\). 已知流出的水中该杂质的质量分数与 \(a\)，\(b\) 的乘积 \(ab\) 成反比. 现有制箱材料 \(60\text{ 平方米}\). 问当 \(a\)，\(b\) 各为多少米时，经沉淀后流出的水中该杂质的质量分数最小. （\(A\)，\(B\) 孔的面积忽略不计）

\bitmapfigure[max width=6.2cm]{img/1998/national_science/q22_fig1.png}
\end{problem}

\begin{solution}
设流出的水中杂质的质量分数为 \(y\). 由题意，存在常数 \(k>0\)，使
\[
y=\frac{k}{ab}
\]
因此只需在约束条件下使 \(ab\) 最大. 无盖箱体所需材料面积为底面、两个长侧面和两个端面的面积之和，所以
\[
2a+2ab+4b=60
\]
即 \(a+ab+2b=30\)，且 \(a>0\)，\(b>0\).
由基本不等式，
\[
a+2b\geq2\sqrt{2ab}
\]
从而令 \(u=\sqrt{ab}>0\)，有
\[
u^{2}+2\sqrt{2}u\leq30
\]
即 \((u+\sqrt{2})^{2}\leq32\).
所以 \(u\leq3\sqrt{2}\)，即 \(ab\leq18\). 等号成立当且仅当 \(a=2b\). 结合 \(ab=18\)，得
\[
2b^{2}=18
\]
由于 \(b>0\)，所以 \(b=3\)，\(a=6\).

因此当 \(a=6\text{ 米}\)，\(b=3\text{ 米}\) 时，\(ab\) 取得最大值，质量分数 \(y\) 取得最小值.
\end{solution}

\begin{problem}
已知斜三棱柱 \(ABC - A_{1}B_{1}C_{1}\) 的侧面 \(A_{1}ACC_{1}\) 与底面 \(ABC\) 垂直，\(\angle ABC = 90^{\circ}\)，\(BC = 2\)，\(AC = 2\sqrt{3}\)，且 \(AA_{1} \perp A_{1}C\)，\(AA_{1} = A_{1}C\).

\begin{enumerate}
\item 求侧棱 \(A_{1}A\) 与底面 \(ABC\) 所成角的大小；
\item 求侧面 \(A_{1}ABB_{1}\) 与底面 \(ABC\) 所成二面角的大小；
\item 求顶点 \(C\) 到侧面 \(A_{1}ABB_{1}\) 的距离.
\end{enumerate}

\bitmapfigure[max width=6.2cm]{img/1998/national_science/q23_fig1.png}
\end{problem}

\begin{solution}
\begin{enumerate}
\item 在平面 \(A_{1}ACC_{1}\) 内作 \(A_{1}D\perp AC\)，垂足为 \(D\). 因为侧面 \(A_{1}ACC_{1}\) 与底面 \(ABC\) 垂直，所以 \(A_{1}D\perp\) 平面 \(ABC\). 又 \(\triangle AA_{1}C\) 在 \(A_{1}\) 处为等腰直角三角形，故 \(AC\) 是其斜边，\(D\) 是 \(AC\) 的中点，并且
\[
AD=A_{1}D=\frac{AC}{2}=\sqrt{3}
\]
于是 \(\angle A_{1}AD\) 是侧棱 \(A_{1}A\) 与底面 \(ABC\) 所成的角，且
\[
\tan\angle A_{1}AD=\frac{A_{1}D}{AD}=1
\]
所以所成角为 \(45^{\circ}\).

\item 过 \(D\) 作 \(DE\perp AB\)，垂足为 \(E\). 因为 \(AB\perp BC\)，所以 \(DE\parallel BC\). 又 \(D\) 是 \(AC\) 的中点，故 \(E\) 是 \(AB\) 的中点，从而
\[
DE=\frac{BC}{2}=1
\]
直线 \(A_{1}E\) 在侧面 \(A_{1}ABB_{1}\) 内，且 \(A_{1}D\perp AB\)、\(DE\perp AB\)，所以 \(A_{1}E\perp AB\). 因而 \(\angle A_{1}ED\) 是侧面 \(A_{1}ABB_{1}\) 与底面 \(ABC\) 所成二面角的平面角. 在直角三角形 \(A_{1}DE\) 中，
\[
\tan\angle A_{1}ED=\frac{A_{1}D}{DE}=\sqrt{3}
\]
所以所成二面角为 \(60^{\circ}\).

\item 设点 \(C\) 到平面 \(A_{1}ABB_{1}\) 的距离为 \(h\). 由勾股定理，
\[
A_{1}E=\sqrt{A_{1}D^{2}+DE^{2}}=\sqrt{3+1}=2
\]
于是
\(S_{\triangle A_{1}AB}=\frac12 AB\cdot A_{1}E=AB\)，\(S_{\triangle ABC}=\frac12 AB\cdot BC=AB\)
四点 \(C\)，\(A_{1}\)，\(A\)，\(B\) 确定的三棱锥无论取 \(\triangle A_{1}AB\) 还是 \(\triangle ABC\) 为底面，体积相同，因此
\[
\frac13S_{\triangle A_{1}AB}h=\frac13S_{\triangle ABC}A_{1}D
\]
由两个底面积相等，得 \(h=A_{1}D=\sqrt{3}\). 所以点 \(C\) 到侧面 \(A_{1}ABB_{1}\) 的距离为 \(\sqrt{3}\).
\end{enumerate}
\end{solution}

\begin{problem}
设曲线 \(C\) 的方程是 \(y = x^{3} - x\)，将 \(C\) 沿 \(x\) 轴，\(y\) 轴正向分别平行移动 \(t\)，\(s\) 单位长度后得曲线 \(C_{1}\).

\begin{enumerate}
\item 写出曲线 \(C_{1}\) 的方程；
\item 证明曲线 \(C\) 与 \(C_{1}\) 关于点 \(A\left(\frac{t}{2}, \frac{s}{2}\right)\) 对称；
\item 如果曲线 \(C\) 与 \(C_{1}\) 有且仅有一个公共点，证明 \(s = \frac{t^{3}}{4} - t\) 且 \(t \neq 0\).
\end{enumerate}
\end{problem}

\begin{solution}
\begin{enumerate}
\item 平移后横坐标增加 \(t\)，纵坐标增加 \(s\). 因此曲线 \(C_{1}\) 的方程为
\[
y=(x-t)^{3}-(x-t)+s
\]

\item 在曲线 \(C\) 上任取一点 \(P=(x,y)\)，则 \(y=x^{3}-x\). 点 \(P\) 关于 \(A\left(\frac{t}{2},\frac{s}{2}\right)\) 的对称点为
\[
P'=(t-x,s-y)
\]
把 \(P'\) 的横坐标代入 \(C_{1}\) 的方程右端，得
\[
(t-x-t)^{3}-(t-x-t)+s=(-x)^{3}+x+s=s-y
\]
所以 \(P'\) 在曲线 \(C_{1}\) 上. 反过来，若 \(Q=(u,v)\) 在曲线 \(C_{1}\) 上，则 \(v=(u-t)^{3}-(u-t)+s\)，其关于点 \(A\) 的对称点为 \(Q'=(t-u,s-v)\)，并且 \(s-v=(t-u)^{3}-(t-u)\)，所以 \(Q'\) 在曲线 \(C\) 上. 因而曲线 \(C\) 与 \(C_{1}\) 关于点 \(A\) 对称.

\item 两曲线公共点的横坐标必须满足
\[
x^{3}-x=(x-t)^{3}-(x-t)+s
\]
整理得
\[
3tx^{2}-3t^{2}x+t^{3}-t-s=0
\]
若 \(t=0\)，则当 \(s=0\) 时两曲线重合，公共点有无穷多个；当 \(s\ne0\) 时两曲线没有公共点，所以恰有一个公共点时必有 \(t\ne0\). 此时上式是一元二次方程，恰有一个公共点等价于它有一个重根，故
\[
\Delta=(-3t^{2})^{2}-4(3t)(t^{3}-t-s)=-3t^{4}+12t^{2}+12ts=3t(-t^{3}+4t+4s)=0
\]
结合 \(t\ne0\)，得
\[
s=\frac{t^{3}}{4}-t
\]
\end{enumerate}
\end{solution}

\begin{problem}
已知数列 \(\{b_{n}\}\) 是等差数列，\(b_{1} = 1\)，\(b_{1} + b_{2} + \cdots + b_{10} = 145\).

\begin{enumerate}
\item 求数列 \(\{b_{n}\}\) 的通项 \(b_n\)；
\item 设数列 \(\{a_{n}\}\) 的通项 \(a_n = \log_a \left(1 + \frac{1}{b_n}\right)\)（其中 \(a > 0\)，且 \(a \neq 1\)），记 \(S_n\) 是数列 \(\{a_{n}\}\) 的前 \(n\) 项和. 试比较 \(S_n\) 与 \(\frac{1}{3}\log_a b_n + 1\) 的大小，并证明你的结论.
\end{enumerate}
\end{problem}

\begin{solution}
\begin{enumerate}
\item 设数列 \(\{b_{n}\}\) 的公差为 \(d\). 由等差数列前 \(10\) 项和公式，
\[
145=\frac{10}{2}\left(2b_{1}+9d\right)=5(2+9d)
\]
解得 \(d=3\)，所以
\[
b_{n}=b_{1}+(n-1)d=3n-2
\]

\item 由 \(b_{n}=3n-2>0\)，有
\[
S_{n}=\sum_{k=1}^{n}\log_{a}\left(1+\frac{1}{3k-2}\right)=\log_{a}P_{n}
\]
其中
\[
P_{n}=\prod_{k=1}^{n}\frac{3k-1}{3k-2}
\]
而 \(\frac{1}{3}\log_{a}b_{n}+1=\log_{a}\left(b_{n}^{1/3}\right)+\log_{a}a=\log_{a}\left(ab_{n}^{1/3}\right)\). 令 \(R_{n}=\frac{P_{n}}{b_{n}^{1/3}}\)，则 \(S_{n}-\left(\frac{1}{3}\log_{a}b_{n}+1\right)=\log_{a}\frac{R_{n}}{a}\).

下面证明 \(R_{n}>1\). 实际上证明更强的不等式 \(P_{n}^{3}>3n+1\). 当 \(n=1\) 时，\(P_{1}=2\)，故 \(P_{1}^{3}=8>4\). 假设 \(P_{k}^{3}>3k+1\)，则
\[
P_{k+1}^{3}=P_{k}^{3}\left(\frac{3k+2}{3k+1}\right)^{3}>\frac{(3k+2)^{3}}{(3k+1)^{2}}>3k+4
\]
其中最后一个不等式由
\[
(3k+2)^{3}-(3k+4)(3k+1)^{2}=9k+4>0
\]
得到. 因此 \(P_{n}^{3}>3n+1>3n-2=b_{n}\)，从而 \(R_{n}>1\).

所以，当 \(0<a<1\) 时，\(\frac{R_{n}}{a}>1\)，而 \(\log_{a}x\) 在 \(x>0\) 上单调递减，故
\[
S_{n}<\frac{1}{3}\log_{a}b_{n}+1
\]
当 \(a>1\) 时，\(\log_{a}x\) 单调递增，比较结果须按 \(a\) 与 \(R_{n}\) 的大小分为三种情形：
\[
\begin{cases}
S_{n}>\frac{1}{3}\log_{a}b_{n}+1,&1<a<R_{n},\\
S_{n}=\frac{1}{3}\log_{a}b_{n}+1,&a=R_{n},\\
S_{n}<\frac{1}{3}\log_{a}b_{n}+1,&a>R_{n}
\end{cases}
\]
其中
\[
R_{n}=\frac{1}{(3n-2)^{1/3}}\prod_{k=1}^{n}\frac{3k-1}{3k-2}
\]
因此，当 \(a>1\) 时，比较结果由 \(a\) 与上述精确阈值 \(R_{n}\) 的大小决定；不能仅由 \(a>1\) 得到对所有 \(n\) 都相同的大小关系. 综上，这三种情形与 \(0<a<1\) 时的结论构成完整比较结果.
\end{enumerate}
\end{solution}
